\scanpage{59}
\noindent\hfill\unsure{atkin page 3}\mspage{3}\par
\begingroup
\csname tagsleft@true\endcsname

We now define the Hecke operator $\handwritten{\underline}{T_p}$ for prime $\handwritten{\underline}{p}$ acting on a modular form $\handwritten{\underline}{\deleted{\unsure{E|F}}\handwritten{F}(z)}$ of dimension $\handwritten{\underline}{k}$ by
\[
T_p\,F(z)=p^{-1}\handwritten{\sum_{\underline{\lambda}=0}^{p-1}F\left(\frac{z+\underline{\lambda}}{p}\right)}
+p^{-k-1}F(pz),
\]
\margin{\handwritten{Lambda}}
or equivalently in terms of the FS:
\[
T_p\,F(x)=\handwritten{\sum_n}\handwritten{\Bigl\{}a(np)+\unsure{.}p^{-k-1}\,a(n/p)\handwritten{\Bigr\}}\,x^n\qquad ,
\]
where $\handwritten{\underline}{a}(\handwritten{\underline}{\alpha})$ is defined as zero if $\handwritten{\underline}{\alpha}$ is non-integral or $\handwritten{\underline}{\alpha}<\handwritten{\underline}{N}$.\margin{\handwritten{alpha}} It is then a matter of formal algebra to establish the two following properties:
\[
\text{\handwritten{\uline}{$T_p$ maps $C$ into $C$,and $C^{\handwritten{0}}$ into $C^{\handwritten{0}}$}}\ ,
\tag*{(H1)}
\]
\[
\text{\handwritten{\uline}{All the $T_p$ commute.}}
\tag*{(H2)}
\]
The further development of the classical theory now requires Petersson's scalar(inner)product
\[
(F,G)=\handwritten{\iint_{D}}\,F(z)\,\handwritten{\overline}{G(z)}\,\handwritten{\underline{\eta}^{-k-2}\,d\underline{\xi}\,d\underline{\eta}}
\qquad (z=\handwritten{\underline{\xi}+i\underline{\eta}})\ ,
\]
\margin{\handwritten{gsi\\ eta}}
which co\handwritten{n}verges provided $\handwritten{\underline}{F,G}\handwritten{\in}\handwritten{\underline}{C}^{\handwritten{+}}$ and at least one of $\handwritten{\underline}{F,G}\handwritten{\in}\handwritten{\underline}{C}^{\handwritten{0}}$, and makes $\handwritten{\underline}{C}^{\handwritten{0}}$ into a Hilbert space. We then have
\[
\text{\handwritten{\uline}{$T_p$ is an hermitian operator on $C^{\handwritten{0}}$,i.e.}}, 
\tag*{(H3)}
\]
\[
(T_pF,G)=(F,T_pG)\ ,
\]
and we have already stated
\margin{\handwritten{\tikz[x=.6em,y=.6em]{\draw[very thin] (.3,2.1) .. controls (.95,1.7) and (.9,1.4) .. (.17,1.51) .. controls (-.5,1.72) and (-.27,1.12) .. (.25,1.33) .. controls (.75,1.61) and (.5,.52) .. (1,-.6);}}}
\[
\text{\handwritten{\uline}{The Hilbert space $C^{\handwritten{0}}$ is finite-dimensional}.}
\tag*{(H4)}
\]
The properties (H1) through (H4) suffice to establish the following

\handwritten{\uline}{Theorem} 1. (\handwritten{\uline}{Hecke-Petersson}).

\handwritten{\uline}{The space of cuspforms of dimension $k$ on the full modular group has a basis $F_i(z)$ of forms which are simultaneous eigenforms}*\margin{\handwritten{*footnote\\ on cover\\ page}} \handwritten{\uline}{of all the operators $T_p$} . \handwritten{\uline}{In terms of the Fourier series} $\handwritten{\sum\limits_{n=1}^{\infty}}\,a(n)\,x^{n}$ \handwritten{\uline}{of a typical} $F_i(z)$ \handwritten{\uline}{normalized by} $a(1)=1$ \handwritten{\uline}{this is equivalent to the ``multiplicative'' relation}
\csname tagsleft@false\endcsname
\[
\handwritten{\underline}{a(np)-a(n)a(p)+p^{-k-1}a(n/p)=0}\quad
(\text{\handwritten{\uline}{all prime} }p,\text{\handwritten{\uline}{all integral} }n)\ .
\tag*{( 1 )\,\handwritten{\#}}
\]
From (1) and the values of $\handwritten{\underline}{a(p)}$ for prime $\handwritten{\underline}{p}$ the values of $\handwritten{\underline}{a(n)}$ are completely determined; in particular $\handwritten{\underline}{a(mn)=a(m)a(n)}$ if $\handwritten{\underline}{(m,n)=1}$. For this reason we have not defined the operators $\handwritten{\underline}{T_m}$ for composite $\handwritten{\underline}{m}$ which form a ring,but merely the \deleted{group}\inserted{\handwritten{\unsure{multiplicative|multiplicatif|multiplicat.} \uline{semi} group!}} of $\handwritten{\underline}{T_p}$ ,since the number theory(if not the algebra) is completely described by the equation (1). One can infer Theorem 1 in any specific case with a much weaker form of(H3),and indeed without the scalar product; it suffices to have (H1),(H2),(H4),together with
\csname tagsleft@true\endcsname
\[
\begin{aligned}
&\text{\handwritten{\uline}{At least one} } T_p\ \text{(\handwritten{\uline}{considered as a matrix acting on some basis}}\\
&\text{\unsure{for|on|in} } C^{\handwritten{0}}\text{ ) has distinct eigenvalues.}
\end{aligned}
\tag*{(H3\unsure{*|$\cdot$})}
\]
\endgroup

% archive sheet number: 58
